\documentclass[11pt,compress,t,notes=noshow, xcolor=table]{beamer}

\input{../../style/preamble}
\input{../../latex-math/basic-math}
\input{../../latex-math/basic-ml}

\title{Optimization in Machine Learning}

\begin{document}

\titlemeta{
First order methods
}{
In practice
}{
figure_man/linesearch.png
}{
\item Optimizers are pre-implemented in major frameworks
\item Minimal training loop in PyTorch
\item Hyperparameters and sensible defaults
}

\begin{framei}{Optimizers are already implemented}
\item All methods we saw (GD, SGD, $\dots$, Adam) are pre-implemented in major ML frameworks
\item You rarely implement them yourself -- just configure and call
\item This reflects how generic and effective they are
\spacer
\item[Python] PyTorch (\texttt{torch.optim}), JAX (\texttt{optax}), TensorFlow / Keras (\texttt{keras.optimizers})
\item[R] available via \texttt{\{torch\}} and \texttt{\{keras\}}; base \texttt{optim()} instead provides quasi-Newton methods (see Chapter 5)
\end{framei}


\begin{frame}[fragile]{Minimal example: Adam in PyTorch}

\begin{verbatim}
# model: a torch.nn.Module; loss_fn: e.g. MSELoss

opt = torch.optim.Adam(model.parameters(), lr=1e-3)

for x, y in dataloader:          # mini-batches
    opt.zero_grad()              # reset gradients
    loss = loss_fn(model(x), y)  # forward pass
    loss.backward()              # backprop -> .grad
    opt.step()                   # Adam update
\end{verbatim}

\begin{itemize}
\item The same four calls drive every optimizer: swap \texttt{Adam} for \texttt{SGD}, \texttt{AdaGrad}, \ldots
\end{itemize}
\end{frame}


\begin{framei}{Hyperparameters: defaults usually work}
\item These methods introduce hyperparameters: step size $\alpha$, stability constant $\beta$, decay rates $\rho_1, \rho_2$, $\dots$
\item \textbf{But:} Defaults are usually strong -- e.g. Adam with $\alpha = 0.001$, $\rho_1 = 0.9$, $\rho_2 = 0.999$, $\beta = 10^{-8}$ works across many tasks
\item $\Rightarrow$ Can be applied effectively without extensive tuning (although tuning can pay off)
\spacer
\item Main knob worth tuning: the \textbf{step size} (and its schedule)
\end{framei}

\endlecture
\end{document}
